\section{CMIN-SR Diagnostic Framework}

\subsection{Overview}

CMIN-SR is designed as a constrained diagnostic pipeline rather than as a
free-form spectrum regressor. The pipeline separates two tasks that are often
mixed together. The first task is empirical spectrum estimation: from a finite
signal, obtain a usable \(\widehat{\zeta}(q)\). The second task is spectral
interpretation: decide how this empirical spectrum relates to monofractal and
MRW-family explanations. This separation matters because a clean spectrum and
a noisy spectrum can lead to very different conclusions even when the
downstream interpretation rule is the same.

The framework reports diagnostic scores rather than mechanism probabilities.
The main scores are:
\begin{itemize}
  \item \(p_{\mathrm{scaling}}\): evidence for stable scaling over the tested
  scales;
  \item \(p_{\mathrm{curved}}\): evidence for nonlinear curvature beyond a
  linear monofractal spectrum;
  \item \(p_{\mathrm{mono}}\): compatibility with a monofractal explanation;
  \item \(p_{\mathrm{MRW}}\): compatibility with a curved MRW-family
  projection;
  \item \(p_{\mathrm{boundary}}\): evidence that the spectrum lies near the
  monofractal/MRW boundary.
\end{itemize}
These scores should be read jointly with residuals and instability warnings.
For example, high \(p_{\mathrm{scaling}}\) alone is not evidence of
multifractality, and high MRW compatibility should not be interpreted without
checking the monofractal residual. The notation \(p\) denotes a normalized
diagnostic score produced by the model; it should not be read as a statistically
validated posterior probability.

\subsection{How MRW samples are simulated}

The synthetic MRW samples used in the project are generated in a way that
matches the analytic projection family. A log-volatility modulation
\(\omega_t\) is sampled as a stationary Gaussian field with logarithmic
covariance,
\begin{equation}
  \mathrm{Cov}(\omega_t,\omega_{t+k})
  =
  \lambda^2 \log\frac{L}{k+\epsilon}
  \quad \text{for } k<L,
  \label{eq:log-cov}
\end{equation}
and zero covariance beyond the integral scale \(L\). The implementation samples
this Gaussian field using a circulant embedding of the Toeplitz covariance
matrix. The field is shifted by \(-\mathrm{Var}(\omega_t)/2\), which keeps the
lognormal modulation centered.

For the fGn-based MRW variant used in the CMIN-SR process generator, fractional
Gaussian noise \(\epsilon_t^{(H)}\) is sampled with Hurst exponent \(H\), and
the increment sequence is
\begin{equation}
  \Delta X_t = \sigma \exp(\omega_t)\epsilon_t^{(H)}.
  \label{eq:mrw-increments}
\end{equation}
The path \(X_t\) is obtained by cumulative summation. In most diagnostic
experiments the model receives the standardized increments
\(\Delta X_t\), not the integrated path. The theoretical scaling curve attached
to these samples is Eq.~\eqref{eq:mrw-zeta}. Low-\(\lambda^2\) MRW samples
therefore act as boundary cases: their spectra are close to linear even though
they are generated from the MRW family.

\begin{figure}[t]
  \centering
  \includegraphics[width=0.96\linewidth]{fig9_mrw_simulation_example.pdf}
  \caption{Example MRW simulation used for controlled experiments. The
  log-volatility modulation creates intermittent increments, while the analytic
  spectrum shows how increasing \(\lambda^2\) bends the otherwise linear
  scaling law. The data used to generate this figure are saved with the paper
  assets for reproducibility.}
  \label{fig:mrw-simulation}
\end{figure}

\subsection{Analytic spectrum-space calibration}

To isolate the interpretation problem from the raw estimation problem, the
project first trains and evaluates a small spectral-geometry calibrator on
analytic and semi-analytic \(\zeta(q)\) curves. The analytic dataset includes:
linear monofractal spectra, low-curvature boundary MRW spectra, strongly
curved MRW spectra, noisy monofractal spectra, noisy MRW spectra, heavy-tail
distorted spectra, regime-like apparent spectra, and ambiguous mild-curvature
spectra. Each sample contains the input \(\zeta(q)\), the monofractal and MRW
projection curves, residual norms, MRW-vs-monofractal gain, curvature score,
linearity score, boundary score, and instability targets.

The calibrator itself is intentionally small: it is a multilayer perceptron
over spectrum-space features, not a large raw time-series backbone. Its role is
to learn the geometry of the diagnostic space. This design allows us to ask a
clean question: if the spectrum is correct, can the diagnostic layer separate
linear, boundary, curved, and unstable cases? The answer is positive in the
analytic experiments, which motivates the later failure analysis on raw finite
samples.

\subsection{Raw zeta alignment and curvature preservation}

The project also studies whether a raw encoder can learn a cleaner
\(\widehat{\zeta}(q)\) from finite time series. Raw zeta alignment supervises
the empirical spectrum against analytic targets for MRW, fGn, and Gaussian
controls, using masks for heavy-tailed cases where high-order moments are not
reliable. This stage successfully reduces false curvature for fGn and Gaussian
controls, but it can also over-linearize genuine MRW curvature. A later
curvature-preserving alignment stage adds stronger band-specific curvature
losses for medium/high-\(\lambda^2\) MRW and uses third-difference smoothness
to discourage jagged spectra without directly flattening quadratic curvature.

These stages are important as ablations, but they are not presented as a final
solution to short-window MRW parameter recovery. Their main value is
diagnostic: they show that suppressing fake curvature and preserving real MRW
curvature are competing requirements under finite samples.

\subsection{Analytic spectrum geometry}

Figure~\ref{fig:spectrum-shapes} illustrates the analytic geometry implied by
the projection families. When \(\lambda^2=0\), the MRW curve reduces to the
linear monofractal spectrum. When \(\lambda^2\) is small, the curve is close to
linear and should be treated as a boundary case. When \(\lambda^2\) is larger,
the curvature becomes visible in \(\zeta(q)\), in the slope
\(\alpha(q)\), and in the corresponding \(f(\alpha)\) geometry.

\begin{figure}[t]
  \centering
  \includegraphics[width=0.96\linewidth]{fig7_multifractal_spectrum_shapes.pdf}
  \caption{Analytic monofractal and MRW spectrum geometry. The figure is not an
  empirical reconstruction result; it visualizes the analytic decoder used to
  interpret monofractal, boundary, and curved MRW-family spectra.}
  \label{fig:spectrum-shapes}
\end{figure}

\subsection{Spectral geometry diagnostics}

The final diagnostic vector combines projection and stability quantities:
\[
  \left(
  r_{\mathrm{MRW}},
  r_{\mathrm{mono}},
  \widehat{\lambda}^2_{\mathrm{proj}},
  \mathrm{gain}_{\mathrm{MRW-vs-mono}},
  \mathrm{curvature},
  \mathrm{linearity},
  \mathrm{boundary},
  \mathrm{tail/instability}
  \right).
\]
The diagnostic scores are calibrated from these features. The curvature
diagnostics and \(\widehat{\lambda}^2_{\mathrm{proj}}\) are related but not
identical: the former summarize the observed geometry of
\(\widehat{\zeta}(q)\), whereas the latter is the coordinate obtained by
projecting that geometry onto the MRW family. These quantities should not be
reported in isolation. A useful reading is that \(p_{\mathrm{MRW}}\) is
meaningful only when scaling is stable, curvature is present, the MRW residual
is lower than the monofractal residual, and instability warnings are not
dominant.
This conditional interpretation is the central design principle of CMIN-SR.

\begin{figure}[t]
  \centering
  \includegraphics[width=0.72\linewidth]{fig2_spectral_geometry_calibration.pdf}
  \caption{Analytic spectral geometry calibration. Clean monofractal,
  boundary-MRW, and curved-MRW spectra occupy distinct regions of calibrated
  spectral diagnostic space.}
  \label{fig:spectral-geometry}
\end{figure}
